\section{Inference: derivations and implementation notes}
\label{app:inference}

This appendix supplies the derivations \S\ref{sec:belief} states without proof
and the implementation detail a reimplementer needs. It also carries the
exactness audit in full, because the headline result of \S\ref{sec:belief} is a
comparison against an object claimed to be exact, and a claim of that kind is
worth exactly what its audit is worth.

The source files are \filepath{engine/src/belief.hpp} (constraint bookkeeping and
the approximate path), \filepath{engine/src/blockdp.hpp} (the exact count law),
and \filepath{engine/src/oracle.hpp} (the brute-force enumerator). Notation is
that of \S\ref{sec:belief-state}.

\subsection{The constraint system}
\label{app:inf-constraints}

The observer's state is a \texttt{Knowledge} object: a resolved-owner array, a
per-card surviving-support bitmask $M_c$, the public hand counts, per-seat and
per-half-suit ask and miss tallies, and a vector of live ask-legality
certificates. (C1)--(C3) are maintained incrementally as bit operations on the
support masks, and a support that collapses to a single seat resolves the card,
which in turn removes any certificate the resolution satisfies.

(C4) is derived rather than stored. Capacities are recomputed from the public
hand counts by \eqref{eq:capacity} after every event, and a seat with no spare
capacity is removed from the support of every unresolved card, which can cascade.
This is the only place the two constraint families interact, and it is why a
declaration --- correct or incorrect --- needs no special handling anywhere in
the system. A declaration removes six cards from play and changes the public
counts of the seats that held them; the six cards leave $U$ with their holders
known, the half-suit leaves $\mathcal{H}$, and each $h_p$ in
\eqref{eq:capacity} drops by the number of those cards that seat held. The same
identity that absorbs a successful ask absorbs a declaration. An incorrect
declaration is not a special case either: it reveals the true holders of all six
cards, which is strictly more information than a correct one gives the opposing
team, and it arrives through the same channel.

(C5) is stored as a list. Each certificate is a pair --- the asking seat and the
bitmask $D$ of \eqref{eq:disj} --- and a certificate whose card set collapses to
a single card resolves that card immediately rather than being stored. Two
properties of the store matter downstream. First, every certificate is confined
to one half-suit, which is what \S\ref{app:inf-blocks} exploits. Second, the
store admits duplicates: a repeated ask in the same half-suit under the same
support produces an identical certificate, and \texttt{onEvent} appended it
without checking. The de-duplication repair of \S\ref{sec:engine-repairs} is
opt-in for a reason that belongs in this appendix rather than in the main text:
the approximate conditioning step of \S\ref{app:inf-sinkhorn} applies each stored
certificate once, so two copies of one certificate condition twice and the
duplicate is not inert. De-duplicating is therefore a change of posterior, not a
change of representation, and it is off by default so that the \vfour{} and
\vfive{} numbers this report compares against are unchanged.

\subsection{The capacity-vector count law}
\label{app:inf-count}

\paragraph{The layer identity.} Order the unresolved cards $c_1,\dots,c_{|U|}$
and let $q$ be the capacity vector. A partial assignment of the first $j$ cards
leaves a vector $n$ of unused capacities, and because every one of those cards
consumes exactly one unit of capacity from exactly one seat,
\[
  \sum_p n_p \;=\; \sum_p q_p - j \;=\; |U| - j ,
\]
so $j$ is recoverable from $n$ and the layer index need not be stored. That is
what collapses the natural $|U|$-layer table into two single arrays. Each is laid
out as one array of size $\prod_p (q_p+1)$, addressed by the mixed-radix offset
$\mathrm{off}(n) = \sum_p n_p s_p$ with strides $s_p = \prod_{r<p}(q_r+1)$, and a
companion index sorts the reachable states by $\sum_p n_p$ so that a sweep visits
one layer at a time. Count vectors are packed one byte per seat into a single
\verb|uint64_t|, which makes the dominance test $n \ge t$ that guards a block
transition one branch-free word operation.

\paragraph{The bound.} The observer's own cards are resolved by (C1), so $q_i =
0$ and the table is a product over the five other seats. Also $|U| \le 45$: at
the deal the 54 cards lie in nine live half-suits, the observer holds nine and
those are resolved, and thereafter a card leaves $U$ when its holder is revealed
or its half-suit is declared and never re-enters. Subject to $\sum_{p\ne i} q_p =
|U|$, the product $\prod_{p \ne i}(q_p+1)$ is maximised at equal capacities, by
the arithmetic--geometric mean inequality, so
\begin{equation}
  \prod_{p\ne i}(q_p+1) \;\le\; \Bigl(\frac{\sum_{p \ne i}(q_p+1)}{5}\Bigr)^{5}
  \;=\; \Bigl(\frac{|U|}{5}+1\Bigr)^{5} \;\le\; 10^{5}.
  \label{eq:bound}
\end{equation}
Each state has at most six outgoing transitions, so a full forward--backward pass
costs at most $6 \times 10^5$ multiply--adds regardless of the position and the
two arrays together occupy under two megabytes. The implementation refuses to
build above a fixed state cap and reports the failure to the caller rather than
degrading silently.

\paragraph{The per-card marginal.} The marginal is the contraction of the two
tables at card $j$,
\begin{equation}
  \mu_{c_j, p}
    \;=\; \frac{\mathbf{1}[\,p \in M_{c_j}\,]}{Z}
      \sum_{\substack{n \,:\, |n| = |U| - j + 1 \\ n_p \ge 1}}
        F_{j-1}(n)\; B_j(n - e_p),
  \label{eq:marg}
\end{equation}
which is exact: every consistent assignment is counted once, by splitting it at
card $j$ into a prefix counted by $F_{j-1}$ and a suffix counted by $B_j$. Both
tables hold integer-valued counts in double precision, so the arithmetic is exact
while the counts stay below $2^{53}$, which \eqref{eq:bound} guarantees on every
reachable state.

\paragraph{The direct sampler.} The backward table doubles as a sampler. Having
placed $c_1,\dots,c_{j-1}$ and arrived at remaining capacities $n$, card $c_j$
goes to seat $p$ with probability proportional to $B_j(n - e_p)$, restricted to
$p \in M_{c_j}$ with $n_p \ge 1$. A seat receives positive weight only when the
suffix can still be completed, so every prefix is extendable and no draw is
discarded: the procedure is rejection-free for (C1)--(C4) and uniform over the
assignments satisfying them, because the weights are path counts. When a live
(C5) certificate is present the sampler is unchanged and the completed assignment
is tested against the literal disjunction \eqref{eq:disj}, drawing again on
failure, so the combined procedure is exact but not rejection-free. The sampler
is off the deployed decision path and is used as ground truth in the
cross-checks of \S\ref{app:inf-audit}.

\subsection{Half-suit block enumeration}
\label{app:inf-blocks}

$U$ is partitioned by half-suit and each group is classified.

\begin{description}[style=nextline]
  \item[Card group.] No live certificate in that half-suit. The group's cards are
        expanded one at a time by \eqref{eq:fwd}--\eqref{eq:bwd}, branching factor
        $|M_c| \le 5$.
  \item[Block group.] At least one live certificate. The group's assignments are
        enumerated exhaustively by depth-first search over the supports --- at
        most $5^6 = 15{,}625$ candidates for a six-card half-suit --- each tested
        against every certificate of that half-suit, and the survivors aggregated
        by count vector into $g_H(t)$ of \eqref{eq:blocktable}. Alongside
        $g_H(t)$ the engine accumulates the per-card breakdown
        $g_H^{(c\to p)}(t)$ that \eqref{eq:q1} needs, so one enumeration serves
        all three queries.
\end{description}

A group of either kind consumes a known number of cards, so the layer identity
survives and the outer recursion runs over groups: $F_b(n) = \sum_t g_b(t)
F_{b-1}(n+t)$, with $g_b$ the indicator of a single card placement in the
card-group case. Nothing in this treatment is approximate: the enumeration is
exhaustive, the certificate test is the literal disjunction, and the outer
recursion is a counting argument. Inclusion--exclusion over $k$ certificates
would instead cost $2^k$ passes over the whole of $U$, which is the cost the
half-suit locality of (C5) avoids.

\begin{corollary}[Equal probability within a count vector]
\label{cor:equalprob}
Under the uniform posterior of \S\ref{sec:belief-state}, any two allocations of
the same half-suit that satisfy (C1)--(C5) and share a count vector have equal
probability.
\end{corollary}

\begin{proof}
By \eqref{eq:q2} the probability of a surviving allocation $A$ depends on $A$ only
through $t(A)$, since $S_t$ is a function of the count vector and of the tables
outside the block.
\end{proof}

The maximum-probability allocation is therefore decided by the count vector
alone, and the choice among survivors within a count vector is arbitrary. The
caveat is the operative half: an allocation that violates a certificate has
probability zero \emph{even when its count vector is shared by survivors}, so
survival must be checked explicitly and cannot be inferred from the count vector.
The count-vector table records only survivors, which makes the corollary usable
and makes the caveat necessary at the same time. \S\ref{app:inf-latent} records
which of the two consumers of this table actually performs the check.

\subsection{The approximate conditioning step}
\label{app:inf-sinkhorn}

The approximate path maintains a matrix $p_{c,p}$ over unresolved cards and
seats, initialised to the prior weights \eqref{eq:prior} on the support $M_c$ and
zero off it. One Sinkhorn iteration normalises each row to sum to one, then
scales each column $p$ by $q_p / \sum_c p_{c,p}$, enforcing \eqref{eq:capacity}
in expectation. The deployed setting is four outer sweeps of eight such
iterations, with a conditioning step for the certificates applied between
consecutive sweeps and a final row normalisation.

The conditioning step derives as follows. Fix a certificate with asking seat $a$
and card set $D$, and let $E$ be the event $\bigvee_{d\in D}[\sigma(d)=a]$.
Treating $\{p_{d,a}\}_{d \in D}$ as independent Bernoulli indicators,
\[
  \Pr[\neg E] \;=\; \prod_{e \in D} (1 - p_{e,a}),
  \qquad
  \Pr[E] \;=\; 1 - \prod_{e \in D} (1 - p_{e,a}) .
\]
For $d \in D$ the event $\sigma(d) = a$ implies $E$, so $\Pr[\sigma(d) = a \mid
E] = p_{d,a}/\Pr[E]$. For a seat $p \ne a$, the event $\sigma(d) = p$ is
compatible with $E$ only if some other card of $D$ goes to $a$, whose independent
probability is $1 - \prod_{e \in D\setminus\{d\}}(1 - p_{e,a})$. Together,
\begin{equation}
  \Pr[\sigma(d) = a \mid E] = \frac{p_{d,a}}{1 - \prod_{e \in D}(1 - p_{e,a})},
  \qquad
  \Pr[\sigma(d) = p \mid E] = p_{d,p}\,
    \frac{1 - \prod_{e \in D \setminus \{d\}}(1 - p_{e,a})}{1 - \prod_{e \in D}(1 - p_{e,a})} .
  \label{eq:cond}
\end{equation}
Rows are renormalised afterwards and the next sweep restores the capacity
constraints the reweighting disturbed. Two degenerate cases are guarded: when
$\Pr[E]$ underflows the mass is forced onto the certificate rather than dividing
by a near-zero denominator, and when $\Pr[\neg E]$ underflows the certificate is
already implied and the step is skipped.

The independence assumption in \eqref{eq:cond} is false, because the capacity
constraints couple the cards of $D$, so interleaving it with capacity fitting does
not converge to the exact marginals of \eqref{eq:q1}. That is the sense in which
the deployed path is approximate, and it is the only sense: the bookkeeping of
(C1)--(C4) it runs on is exact.

\paragraph{The deployed declaration estimator.} The declaration query
\eqref{eq:q2} asks for a joint probability and the approximate path has only
marginals, so it evaluates the joint by sequential conditioning. For a candidate
allocation assigning cards $c_1,\dots,c_6$ to seats $p_1,\dots,p_6$,
\[
  \hat\alpha(A) \;=\; \prod_{j=1}^{6} \hat\mu^{(j)}_{c_j, p_j},
\]
where $\hat\mu^{(j)}$ is recomputed on the state in which $c_1,\dots,c_{j-1}$
have already been resolved to $p_1,\dots,p_{j-1}$. Each step fixes one card,
decrements that seat's capacity, propagates the support reductions and refits. A
card already resolved contributes one if it matches and zero otherwise, a card
whose support excludes its assigned seat returns zero immediately, and the
product short-circuits once it falls below a floor, since candidates that weak
are rejected regardless. The chain rule is exact --- the joint is the product of
successive conditionals in any fixed order --- so the estimator is exact in
structure and approximate only through the marginal solver. A product of
\emph{unconditioned} marginals would not be: it ignores the capacity coupling
that makes \eqref{eq:q2} necessary in the first place.

\subsection{The policy prior, derived}
\label{app:inf-prior}

Equation \eqref{eq:prior} is a likelihood correction expressed as prior weights.
The full Bayesian posterior over deals weights each deal $x$ by $L(\mathcal{I}
\mid x) = \prod_t \pi_{p_t}(a_t \mid \mathcal{I}_{<t}, x)$, the probability that
the observed actions would have been chosen given $x$. That object is not
available: it requires a model of every opponent's policy and a product over the
whole history. What \eqref{eq:prior} does instead is retain the one sufficient
statistic of the history that any plausible policy makes informative --- how many
of a seat's public asks fell in a given half-suit, against how many fell
elsewhere --- and expose its two coefficients to the same optimiser that fits the
policy. $\theta$ rewards a seat's asks in the card's own half-suit and $\phi$
penalises its asks elsewhere.

The weights enter as the initialisation of the Sinkhorn fit rather than as a
post-hoc multiplication, which is what keeps the capacity constraints exact:
$\theta = \phi = 0$ recovers the uniform initialisation and hence the
policy-agnostic posterior. It also means the prior has no route into the exact
path. \texttt{BlockDP::build} takes the \texttt{Knowledge} object alone; there is
no argument through which a coefficient could reach it, and the exact path never
calls the weight function at all. Any belief-mode ablation in this corpus is
therefore confounded with a prior ablation, and that confound is not a defect of
the experiment but a property of the architecture.

\paragraph{The exact/approximate predictive comparison in full.} The comparison
of \S\ref{sec:belief-exactworse} is run by \texttt{fish v6probe --mode=belief}
and its design is as follows. At every ask decision of the measured team, on the
acting seat's own \texttt{Knowledge} object, each candidate posterior is computed
on the \emph{same} state: the exact count law of \S\ref{sec:belief-count}, and
the approximate path at a list of $(\theta,\phi)$ settings that includes $\theta =
\phi = 0$ and the shipped values. Two quantities are then accumulated per
posterior.

\begin{enumerate}
  \item \textbf{Predictive log loss.} For every unresolved card of the state, the
        true holder is read from the ground-truth deal and $-\log \mu_{c,
        \text{truth}}$ is accumulated, floored to avoid an infinite contribution
        from a posterior that assigns exact zero. The reported figure is the mean
        over cards. The unit is one card, not one decision, so a state with a
        large unresolved pool contributes proportionally more --- which is the
        right weighting, because those are the states inference is for.
  \item \textbf{Argmax hit rate.} The legal asks are enumerated at that state and
        the candidates the acting seat's own knowledge already proves dead are
        discarded. Among the remainder the posterior's own maximiser of
        $\mu_{c,t}$ is selected, and the accumulator records whether the target
        really holds the card. The unit is one decision, and each posterior
        chooses its own ask, so the measure is the quality of the posterior's own
        recommendation rather than of a shared one.
\end{enumerate}

Two properties of this design carry the argument. It is played on identical
states, so no posterior is scored on a trajectory it induced. And it involves no
fitted ask weights on either side --- the argmax is over the posterior itself,
not over the linear utility of \S\ref{sec:mechanisms} --- so nothing in the
comparison depends on which posterior the policy was tuned against. That is what
makes it a verdict on the posteriors rather than a substitution result, and it is
why \S\ref{sec:belief-exactworse} can close a question that the v0.4 study's
belief-substitution ablation could not.

The result is Table~\ref{tab:belief}. The exact count law under a uniform deal
prior scores \vsixBeliefExactNLL{} nats per card and \vsixBeliefExactHit{}\%
argmax hit rate; the approximate path with the prior deleted scores
\vsixBeliefNoPriorNLL{} and \vsixBeliefNoPriorHit{}\%; the shipped prior scores
\vsixBeliefShippedNLL{} and \vsixBeliefShippedHit{}\%. The exact object is last
of the three on both measures. The approximate path beats the exact one even at
$\theta = \phi = 0$, where the two encode the same rules-only information; the
remaining gap to the shipped row is the prior alone, worth
\numThetaPriorHitGain{} points of hit rate and \numThetaPriorNatGain{} nats.

The interpretation is the one \S\ref{sec:belief-exactworse} gives and is repeated
here because the appendix is where a sceptical reader arrives: exactness is
relative to a prior, the uniform deal prior is a modelling assumption and not a
rule, and on this evidence it is a worse assumption than a crude count of who has
asked where.

\subsection{The exactness audit and its coverage}
\label{app:inf-audit}

Agreement between the count law and the block construction is weak evidence,
because they share a derivation. The audit of record is therefore a brute-force
enumerator that factorises nothing: for one observer's constraint state the
unresolved cards are searched by depth-first assignment to seats, candidates are
pruned against the surviving support (carrying (C1)--(C3)) and against the
capacity (carrying (C4)), and (C5) is tested at each complete leaf against the
literal disjunction. Surviving leaves are counted into $Z$, tallied into per-card
ownership counts, and tallied into a per-half-suit table keyed by the packed
assignment. Nothing is factorised, so the count behind every named allocation is
an explicit count of leaves.

That audit was run and reported by the v0.4 study, and this report does not
re-run it \cite{fishbot04}. Its two properties matter here and are restated
rather than re-measured. The deterministic comparisons agree \emph{exactly} ---
not to a tolerance --- because both sides compute integer counts over the same
finite set of assignments and hold them in double precision below $2^{53}$, form
the same ratio from identical operands, and obtain identical doubles. And its
coverage is a subsample: the search is exponential in $|U|$, so a state is
attempted only when the capacity dynamic program reports a small enough number of
consistent deals, and a state that exceeds the budget is counted as skipped and
reported rather than dropped. Skipping is correlated with exactly the property
that makes a state hard --- a large unresolved pool --- so the audit is evidence
about small reachable states and is silent about the high-entropy early game
where inference matters most. A pass is exact on a subsample, and that is all it
is.

Two v0.6 facts attach to this. The first is that the audit's premise did not hold
when it was run. \S\ref{sec:engine-repairs} shows that a \texttt{BlockDP}
instance's tables could be repointed by an unrelated build on the same thread;
the oracle comparison was unaffected because its driver builds one instance and
queries it before anything else can rebuild, which is a property of the driver
rather than of the engine. After the repair it is a property of the engine. The
second is that the exactness audit answers a narrower question than the phrase
suggests: it establishes that the block construction computes the posterior of
(C1)--(C5) correctly, and \S\ref{app:inf-prior} establishes that this posterior
is not the best available one.

The v0.6-era checks that were run are the reachability measurements of
\S\ref{app:inf-latent}.

\subsection{Latent approximations inside functions documented as exact}
\label{app:inf-latent}

Three routines in the inference stack are documented as exact and are not, in
regimes no current caller reaches. Each is recorded here with its reachability
argument, because unreachable-today is a property of the call sites and not of
the code, and every one of these is a trap for the next study rather than an
error in this one. The source is
\filepath{research/v06/notes/R2-inference-stack.md}, \S8.

\paragraph{A greedy branch in the count routine.} The masked counting routine in
\filepath{engine/src/belief.hpp} has a fast path for masks that are all equal and
one for singletons. For masks that are neither --- heterogeneous and
non-singleton --- it falls through to a greedy sequential assignment that picks
the seat with the most remaining capacity, which is an approximation inside a
function whose comment says exact. No current caller reaches it: over
\numCountUniformQueries{} uniform and \numCountGeneralQueries{} general queries
the routine agreed with the exact posterior to
$\numCountAgree{}$. A future caller asking a heterogeneous per-card question ---
``is each of these cards on my team?'' is the natural one --- would get a silent
approximation.

\paragraph{The allocation search does not re-check survival.} The block table's
maximum-probability allocation search materialises an assignment from a count
vector and returns it without re-testing it against the certificates of
\eqref{eq:disj}, whereas the allocation-probability query on the same table does
re-test and says why. The search is on the deployed decision path. It was
measured at \numAllocSurvivalFailures{} failures in \numAllocSurvivalChecks{}
checks, \numAllocSurvivalCert{} of them on states carrying a live certificate,
over \numAllocSurvivalDeals{} deals at seed \numAllocSurvivalSeed{}. That is
unfalsified and it is not proved. Corollary~\ref{cor:equalprob} does not supply
the missing step: it says survivors sharing a count vector are equiprobable, and
says nothing about whether the assignment the search materialises is one of them.

\paragraph{Silent truncation of a group table.} The group enumeration returns a
failure flag when the number of distinct count vectors exceeds
\numMaxEnt{}, but the caller that populates the table ignores the flag and
returns, leaving a \emph{partial} table which the team-ownership query then
normalises as if it were complete. The regime is currently unreachable by
arithmetic: the number of count vectors for six cards over six seats is
$\binom{11}{5} = \numMaxCountVectors{} < \numMaxEnt{}$, and no truncation was
observed in \numAllocSurvivalChecks{} checks. It becomes reachable the moment the
half-suit size or the table capacity moves. Two further caps have the same shape
and the same status: a \numCertCap{}-certificate limit per group, and
\numNodeGuard{}-node guards on the two depth-first searches.
